% =====================================================================
% Section 3: Methodology. Three subsections.
% Before submitting: replace ANONYMOUS-URL-HERE in the release paragraph.
% =====================================================================

\section{Methodology}
\label{sec:method}

\subsection{Task and corpus}
\label{sec:task}
\label{sec:corpus}

An instance is a triple $(v, m, \tau)$: a surveillance video $v$, a source LLM
$m$, and a prompting strategy $\tau$. The source LLM produces a free-form
forensic report $r$ describing what occurs in $v$, and each report carries an
expert reference annotation $g_v$ written independently of any model. The task
is to predict, from $r$ and optionally $g_v$, six binary labels: H1 scene
fabrication, H5 entity fabrication and H6 phantom actors, which invent content
absent from the footage; H3 crime missed, which drops the criminal act; and H2
crime misclassification and H4 severity minimization, which misstate its nature
or gravity. These group onto three axes, \emph{fabrication}, \emph{omission}
and \emph{distortion}. We report per-type AUC and $F_1$, a macro over the six,
a macro restricted to the reliable types (Section~\ref{sec:labels}), and a
binary any-hallucination score.

\paragraph{Provenance.} The reports, taxonomy, panel protocol, per-judge
verdicts and human-validation subset were introduced and released by
prior work \citep{priorwork2026} to measure hallucination prevalence, and \textsc{halt} reuses all of it as raw material. New here are the
held-out-source formulation of the detection task, the reliability tiering, the
four splits, the reference-conditioned and encoder detectors,
and every analysis in Section~\ref{sec:results}.

The inherited corpus is $19{,}361$ reports on 807 UCF-Crime videos
\citep{sultani2018real} with UCA annotations \citep{yuan2024towards} as
reference, written by three frontier multimodal LLMs under eight prompting
strategies. Two of its properties make the held-out-source question askable.
It is fully crossed up to seven failed API calls, so withholding a source LLM
holds
the input fixed by construction. And verbosity varies by a factor of three under an identical
task, 580 mean words per report for GPT against $1{,}802$ for Gemini, the
regularity Section~\ref{sec:results} shows detectors exploiting.

\begin{figure}[!b]
\centering
\includegraphics[width=\textwidth]{fig_pipeline}
\caption{The \textsc{halt} pipeline. The shaded band is inherited; the
remainder is this paper: aggregating the raw verdicts into labels whose
per-type reliability is calibrated against an independently scored subset,
partitioning the corpus into four video-disjoint splits, and evaluating four
detector families against a base-rate control that reads no report text. In
step 2 each cell is one report, rows are source LLMs and columns are
video-by-strategy conditions; existing benchmarks admit only the standard
split.}
\label{fig:pipeline}
\end{figure}


\subsection{Labels and their reliability}
\label{sec:labels}
\label{sec:protocol}
\label{sec:reliability}

The inherited labels come from a panel of the same three models acting as
judges \citep{zheng2023judging}, each report scored by the two that did not
write it, so self-preference bias \citep{panickssery2024llm} is structurally
excluded \citep{priorwork2026}. Three properties of that design matter
here. It yields two verdicts per report under a conjunctive rule, so
alternative rules can be applied post hoc
(Section~\ref{sec:corpusstats}). Judge inputs cap the reference at $1{,}500$
and the report at $3{,}000$ characters, which
Section~\ref{sec:lengthbias} shows shapes the labels. And no model judging its
own output means each source LLM is scored by a distinct judge pair, the
confound of Section~\ref{sec:confound}.

An inherited subset of 130 reports carries independent labels from two human
annotators, members of the research group rather than forensic practitioners,
who worked from text with the panel's labels hidden. From the rater-rater and
panel-consensus agreements we assign each type a reliability tier and exclude
the low tier from the headline macro; the tiering is ours, and
Section~\ref{sec:corpusstats} reports it.

\subsection{Evaluation setup}
\label{sec:splits}
\label{sec:baselines}
\label{sec:release}


Each video contributes 24 reports, so a split over reports would let a detector
memorise the video rather than learn the phenomenon; Figure~\ref{fig:pipeline}
shows the three partitions.
The standard and held-out-strategy splits are therefore video-disjoint. The
held-out-source split is not, since every video is described by every source
LLM; Section~\ref{sec:inversion} reports a disjointness-enforced variant under
which the effect is larger, so the main tables are conservative.

We provide reference detectors in two families. The first is linear: five
logistic regressions with balanced class weights, differing only in features.
\textbf{Surface} uses verbosity, sentence statistics, hedging density, numeric
detail density, type-token ratio, truncation length, and source and strategy
identity indicators; \textbf{surface-noid} drops the identity features;
\textbf{tfidf} uses word 1--2 gram TF-IDF; \textbf{embed} uses sentence
embeddings; and \textbf{grounded} uses report-versus-reference features only,
namely the rate of report tokens absent from the reference, the rate of
reference tokens absent from the report, Jaccard overlap, reference coverage,
length ratio and TF-IDF cosine. Identity features are dropped automatically on
any split where they are constant in training and unseen at test.

The second family is fine-tuned encoders. One encoder per split, six per-type heads plus a
binary head, report and reference as a sentence pair, positive weights matching
the linear baselines. Context length is chosen rather than defaulted: the judge
saw about 1{,}125 tokens, so a detector at 512 sees less than the judge did. We
report \textbf{DeBERTa-v3-base} at its 512-token limit, \textbf{ModernBERT-base}
at 1{,}280, matching the judge's view, and ModernBERT at 2{,}048, exceeding
it.

The corpus, per-judge verdicts, aggregated labels and rubrics were released
with the prior study. \textsc{halt} adds, and releases: the reliability tiers
with supporting $n$, the four splits, reference implementations of every
detector, and a datasheet \citep{gebru2021datasheets}. The artifact is
available at \url{https://anonymous.4open.science/r/HALT-ANON-7FF2/}.
Supporting tables cited below are in Appendices~\ref{app:strategy}
to~\ref{app:rubrics} of this paper.
